\documentclass[12pt]{amsart}

\pdfpagewidth 9in
\pdfpageheight 11in
\topmargin -1.5in
\headheight 0in
\headsep 0in
\textheight 9.5in
\textwidth 9.5in
\oddsidemargin -3in
\evensidemargin 0in
\headheight 77pt
\headsep 0in
\footskip .75in

\usepackage[colorlinks=true,linkcolor=blue,citecolor=blue]{hyperref}

\usepackage{amsmath, amsthm, amssymb,fullpage, xcolor}
\usepackage{bbm}
\usepackage{mathrsfs}
\usepackage{thmtools}
\usepackage{enumerate}
\usepackage{appendix}
\usepackage{mathtools}
\usepackage[all]{xy}


\newcommand{\goe}{G^{N, \bR}}
\newcommand{\gse}{G^{N, \bH} }
\newcommand{\xgoe}{X_N^{\bR}}
\newcommand{\xgse}{X_N^{\bH}}
\newcommand{\pgoe}{\Phi^{\bR}}
\newcommand{\gue}{G^N}

\newcommand{\haaro}{O^{N, \bR}}
\newcommand{\haars}{O^{N, \bH}}

\newcommand{\yo}{Y_N^{\bR}}
\newcommand{\ys}{Y_N^{\bH}}

\newcommand{\infgoe}{\nu^{\mathbb{R}}}
\newcommand{\infgse}{\nu^{\mathbb{H}}}


\newcommand{\po}{\Psi_h^{\bR}}
\newcommand{\pse}{\Psi_h^{\bH}}
\newcommand{\numo}{f_h^{\bR}}
\newcommand{\deno}{g_q^{\bR}}

\newcommand{\bN}{\mathbb{N}}

\newcommand{\wgo}{\wg^O}

\newcommand{\trace}{\mathrm{Tr}}
\newcommand{\tr}{\textnormal{tr}}

\newcommand{\calA}{\mathcal{A}}
\newcommand{\calT}{\mathcal{T}}
\newcommand{\calB}{\mathcal{B}}
\newcommand{\calE}{\mathcal{E}}
\newcommand{\calX}{\mathcal{X}}
\newcommand{\calH}{\mathcal{H}}
\newcommand{\calK}{\mathcal{K}}
\newcommand{\calO}{\mathcal{O}}
\newcommand{\calD}{\mathcal{D}}
\newcommand{\calU}{\mathcal{U}}
\newcommand{\calP}{\mathcal{P}}
\newcommand{\B}{\mathcal{B}}
\newcommand{\calF}{\mathcal{F}}
\newcommand{\jcal}{\mathcal{J}}
\newcommand{\ecal}{\mathcal{E}}
\newcommand{\calM}{\mathcal{M}}
\newcommand{\calv}{\mathcal{V}}

\newcommand{\bZ}{\mathbb{Z}}
\newcommand{\bC}{\mathbb{C}}

\newcommand{\Span}{\mathrm{Span}}

\renewcommand{\P}{\mathbb{P}}
\newcommand{\E}{\mathbb{E}}

\newcommand{\supp}{\mathrm{supp}}

\newcommand{\diag}{\mathrm{diag}}

\newcommand{\Spec}{\mathrm{Spec}}
\newcommand{\spec}{\mathrm{spec}}

\newcommand{\dalpha}{\dot{\alpha}}
\newcommand{\dbeta}{\dot{\beta}}
\newcommand{\dgamma}{\dot{\gamma}}

\newcommand{\les}{\lesssim}

\newcommand{\bS}{\mathbb{S}}

\newcommand{\re}{\mathrm{Re}}
\newcommand{\im}{\mathrm{Im}}

\newcommand{\bw}{w}

\newcommand{\R}{\mathbb{R}}
\newcommand{\C}{\mathbb{C}}
\newcommand{\bR}{\mathbb{R}}

\newcommand{\tv}{\tilde{v}}

\newcommand{\num}[1]{\#(#1)}

\newcommand{\qplus}{\hat{q}}

\newcommand{\rplus}{\mathbb{R}_{>0}}

\newcommand{\op}{\mathrm{op}}

\newcommand{\labs}[1]{\left|#1\right|}

\newcommand{\dif}[1]{\left(\frac{d}{dx}\right)^{#1}}

\newcommand{\cov}{\mathrm{Cov}}

\newcommand{\difff}{\frac{d}{dx}}

\newcommand{\var}{\mathrm{Var}}

\newcommand{\dis}{\mathrm{Dis}}

\newcommand{\an}{A_N}
\newcommand{\bn}{B_N}

\newcommand{\bF}{\mathbb{F}}

\newcommand{\calS}{\mathcal{S}}

\newcommand{\sobrho}[2]{\|#1\|_{C^{#2}[-\rho, \rho]}}
\newcommand{\sobnorm}[2]{\|#1\|_{C^{#2}[-K, K]}}

\newcommand{\abs}[1]{|#1|}

\newcommand{\wg}{\mathrm{Wg}}

\newcommand{\bH}{\mathbb{H}}

\newcommand{\infnorm}[1]{\|#1\|_{C^0[-K, K]}}

\newcommand{\lpl}{K}

\newcommand{\gt}{\mathrm{GT}}

\newcommand{\leb}{\mathrm{Leb}}

\newcommand{\disc}{\mathrm{Disc}}

\newcommand{\vol}{\mathrm{Vol}}

\newcommand{\haar}{\mathrm{Haar}}

\allowdisplaybreaks

\newcommand{\calpha}{\overset{\circ}{\alpha}}
\newcommand{\ceta}{\overset{\circ}{\beta}}

\newcommand{\cu}{\overset{\circ}{u}}
\newcommand{\cw}{\overset{\circ}{w}}

\newcommand{\calV}{\mathcal{V}}
\newcommand{\calW}{\mathcal{W}}

\newcommand{\bone}{\mathbbm{1}}

\newcommand{\rs}{r^{(\sigma)}}

\newcommand{\gdown}{\gamma^{\downarrow}}

\newcommand{\palpha}{E[\alpha|\gamma]}
\newcommand{\pbeta}{E[\beta|\gamma]}

\newcommand{\jaca}{\frac{d\gamma}{d\alpha}}
\newcommand{\jacb}{\frac{d\gamma}{d\beta}}

\newcommand{\jacobian}{\jac}
\newcommand{\jac}{J_{\boxplus_n}}

\newcommand{\roots}{\Omega_{\boxplus_n}}
\newcommand{\rootsi}{\Omega_{\boxplus_n, i}}
\newcommand{\dermap}{ \Omega_{\partial_x}}
\newcommand{\dermapi}{\Omega_{\partial_x, i}}

\newcommand{\score}{\mathscr{J}_n}
\newcommand{\scoreder}{\mathscr{J}_{n-1}}

\newcommand{\jacder}{J_{\partial_x}}

\newcommand{\hessder}{H_{\partial_x}^{(i)}}
\newcommand{\hess}{\mathrm{Hess}}
\newcommand{\hessconv}{H_{\boxplus_n}^{(i)}}

%%%%

\title{The finite free Stam inequality}
\author{Jorge Garza Vargas, Nikhil Srivastava, and  Zack Stier}


\date{\today}

\begin{document}

\maketitle

\numberwithin{equation}{section}

\newtheorem{theorem}{Theorem}[section]
\newtheorem{definition}{Definition}[section]
\newtheorem{conjecture}{Conjecture}[section]
\newtheorem{lemma}{Lemma}[section]
\newtheorem{corollary}{Corollary}[section]
\newtheorem{proposition}{Proposition}[section]
\newtheorem{observation}[theorem]{Observation}
\newtheorem{remark}{Remark}[section]
\newtheorem{problem}[theorem]{Problem}
\newtheorem{assumption}[theorem]{Assumption}
\newtheorem{claim}[theorem]{Claim}
\newtheorem{fact}[theorem]{Fact}
\newtheorem{example}[theorem]{Example}

\maketitle

Let $\boxplus_n$ and $\Phi_n(\cdot)$ be defined as in the problem statement.
In this note we prove the following result, which was conjectured by D. Shlyakhtenko.
\begin{theorem}
\label{thm:stamineq}
    Let $p(x)$ and $q(x)$ be any two monic real-rooted polynomials of degree $n$. Then 
    $$\frac{1}{\Phi_n(p\boxplus_n q)} \geq  \frac{1}{\Phi_n(p)} + \frac{1}{\Phi_n(q)}.$$
\end{theorem}

\section{Notation and preliminaries}
\subsection{Polynomials and the finite free convolution}

Given a polynomial $p(x)$ of degree $n$ we say that $\alpha=(\alpha_1, \dots, \alpha_n)$ is a vector of roots for $p(x)$ if the $\alpha_i$ are the roots of $p(x)$. We will say that $\alpha$ is ordered if $\alpha_1\geq \cdots \geq \alpha_n$.  Recall that for monic polynomials $p(x)$ and $q(x)$, $p(x)\boxplus_n q(x)$ may be expressed as:
\begin{equation}\label{eqn:convroots} p(x)\boxplus_n q(x) = \sum_{\pi\in S_n} \prod_{i=1}^n(x-\alpha_i-\beta_{\pi(i)}),\end{equation}
where $\alpha$ and $\beta$ are vectors of roots for $p(x)$ and $q(x)$, respectively, and $S_n$ is the symmetric group on $n$ elements (see Theorem 2.11 of \cite{marcus2022finite} for a proof). Walsh \cite{walsh1922location} proved that if $p(x)$ and $q(x)$ are real-rooted, then so is $p(x)\boxplus_n q(x)$. Therefore, the finite free convolution induces a map
$$\roots: \bR^{n}\times \bR^n \to \bR^n,$$
where if $\alpha$ and $\beta$ are vectors of roots for $p(x)$ and $q(x)$, then $\roots(\alpha, \beta)$ is defined to be the ordered vector of roots for $p(x)\boxplus_n q(x)$.  

Other than the fact that $\boxplus_n$ preserves real-rootedness, our proof will crucially exploit each of the following well-known properties of the finite free convolution. In what follows we will use $\bone_n$ to denote the all-ones vector of dimension $n$. 
We will use the notation $$m_k(\alpha):= \frac{1}{n}\sum_{i=1}^n \alpha_i^k \quad \text{and} \quad \var(\alpha) := m_2(\alpha)-m_1(\alpha)^2.$$


\begin{proposition}[Properties of $\boxplus_n$]
\label{thm:ffconvolution}
    If $\alpha, \beta\in \bR^n$ and  $\gamma =\roots( \alpha, \beta)$, then:
    \begin{enumerate}[i)]
        \item \label{item:additivity} (Additivity) $m_1(\gamma)=m_1(\alpha)+m_1(\beta)$ and $\var(\gamma) = \var(\alpha)+\var(\beta)$. 
        \item \label{item:translations} (Commutation with translation)   For all $t\in \bR$, $\roots(\alpha+t\bone_n ,\beta) = \gamma +t\bone_n$ and  $\roots(\alpha , \beta + t \bone_n) = \gamma+t\bone_n$.  
    \end{enumerate}
\end{proposition}
\begin{proof}
	(i) Follows from the definition of $p\boxplus_n q$ in terms of the coefficients of $p$ and $q$ and the Newton identities. (ii) Follows from \eqref{eqn:convroots}. 
\end{proof}

 

\subsection{The heat flow and the finite free Fisher information} 
Given a vector of roots $\alpha\in \bR^n$ we will define the its finite free score vector $\score(\alpha) \in (\bR\cup\{\infty\})^n$ as 
$$\score(\alpha) := \left( \sum_{j : j \neq i} \frac{1}{\alpha_i-\alpha_j} \right)_{i=1}^n. $$
Given a real-rooted polynomial $p(x)$ with vector of roots $\alpha$, define its finite free Fisher information as 
$$\Phi_n(p) := \|\score(\alpha)\|^2.$$
The following fact will allow us to write the finite free Fisher information of the polynomial $p(x)$ in terms of the dynamics of its roots under the reverse heat flow. It was shown to us by D. Shlyakhtenko.

\begin{lemma}[Score vectors as derivatives]
\label{lem:score}
	Assume $p(x)$ has simple roots. Let $p_t(x):= \exp\left( -\frac{t}{2} \partial_x^2  \right)p(x)$ and let $\alpha(t) = (\alpha_1(t), \dots, \alpha_n(t))$ be the ordered vector of roots of $p_t(x)$. Then 
     $$\alpha_i'(0) = \sum_{j : j\neq i} \frac{1}{\alpha_i-\alpha_j},$$
    and in particular $\alpha'(0) = \score(\alpha)$. 
\end{lemma}
\begin{proof}
	Since the $\alpha_i(t)$ are continuous in $t$, the roots remain simple in a neighborhood of $t=0$. Implicitly differentiating the expression
	$$p(\alpha_i(t))-tp''(\alpha_i(t))/2+t^2R(\alpha_i(t),t)=0$$
	(where $R(x,t)$ is a polynomial) at $t=0$ one obtains
	$$ \alpha_i'(0)=\frac{1}{2}\frac{p''(\alpha_i)}{p'(\alpha_i)},$$
	which is equal to the advertised expression.
\end{proof}


\section{Proof of Stam's inequality}
We now prove Theorem \ref{thm:stamineq}.  The following Lemma allows us to restrict attention to the case when $p,q$, and $p\boxplus_n q$ all have simple roots.

\begin{lemma}[Approximation by Simple Rooted Polynomials] Let $\epsilon>0$ and define the differential operator $T_\epsilon:=(1-\epsilon \cdot d/dx)^n$. If $p(x)$ is a monic real-rooted polynomial of degree $n$, then 
	\begin{enumerate}[i)]
	\item $(T_\epsilon p)(x)$ is monic and real-rooted of degree $n$ with simple roots. 
	\item $ \Phi_n(T_\epsilon p)\rightarrow \Phi_n(p)$ as $\epsilon\rightarrow 0$. 
	\item  $(T_\epsilon p)\boxplus_n (T_\epsilon q)= T_\epsilon^2 (p\boxplus_n q)$. 
	\end{enumerate}
\end{lemma}
	\begin{proof} (i) was shown in \cite{nuij1968note}. (ii) is because $\Phi_n$ is continuous in the roots of $p$, which are continuous in $\epsilon$. (iii) follows because $\boxplus_n$ commutes with differential operators (see e.g. \cite{marcus2022finite}.)\end{proof}
		Thus, establishing Theorem \ref{thm:stamineq} for the simple case implies the general case by using (iii) above and taking $\epsilon\rightarrow 0$. In what follows, $p(x)$ and $q(x)$ are monic real-rooted polynomials, $\alpha$ and $\beta$ are vectors of roots for $p(x)$ and $q(x)$, $\gamma:= \roots(\alpha, \beta)$, and $\alpha,\beta,\gamma$ all have distinct entries, implying that they are smooth functions of the coefficients of the corresponding polynomials. Let $\jac$ denote the Jacobian of $\roots$ at the point $(\alpha, \beta)$.

Our proof can be separated into three steps. The second step is  the most substantial one and we will defer its detailed discussion to Section \ref{sec:understandingjacob}. 

\textbf{Step 1 (Jacobians and score vectors).} We first note that the following relation between score vectors holds. 


\begin{observation}[Relating score vectors]
\label{obs:scorevecs}
Using the above notation, for any $a, b \geq 0$ 
$$\jac (a\score(\alpha), b \score(\beta))  = (a+b)\score(\gamma). $$
\end{observation}

\begin{proof}
For every $t\geq 0$ let $p_t(x) = \exp(-\tfrac{t}{2} \partial_x^2) p(x)$, let $\alpha(t)$ be the ordered vector of roots of $p_t$, and define $q_t, r_t$ and $\beta(t), \gamma(t)$ in an analogous way. Since the finite free convolution commutes with any differential operator,  it follows that
      $$r_{(a+b)t} = p_{at} \boxplus_n q_{bt}. $$
    Hence  $\gamma((a+b)t) = \roots(\alpha_{at}, \beta_{bt})$ for every $t$. So, if we differentiate this  relation with respect to $t$, using the chain rule for the right-hand side,  we get
      $$(a+b)\gamma'(0) =  \jac \left(\begin{array}{c}
           a \cdot \alpha'(0)  \\ b \cdot\beta'(0)
      \end{array} \right).$$  
    A direct application of  Lemma \ref{lem:score} concludes the proof. 
\end{proof}


\textbf{Step 2 (Understanding the Jacobian).} The substance of our proof lies in understanding  $\jac$. In particular, we will show the following. 

\begin{proposition}
\label{prop:singularvalues}
  If $u, v \in \bR^n$ are orthogonal to $\bone_n$ then
    $$\|\jac (u, v)\|^2 \leq \|u\|^2 + \|v\|^2.$$
\end{proposition}

This proposition will be proven in Section \ref{sec:understandingjacob}, for now we show how it is used. 


\textbf{Step 3 (Proof of Theorem \ref{thm:stamineq} \`a la Blachman).} With Observation \ref{obs:scorevecs} and Proposition \ref{prop:singularvalues} in hand we can conclude the proof using the same argument that Blachman used in \cite{blachman2003convolution}. 

\begin{proof}[Proof of Theorem \ref{thm:stamineq}]
    First note that 
    $$\sum_{i=1}^n \sum_{j : j\neq i} \frac{1}{\alpha_i-\alpha_j}=0,$$
   since each term in the sum appears once with a plus and once with a minus. Therefore $\score(\alpha)$ is orthogonal to $\bone_n$ and, arguing analogously,  $\score(\beta)$ is orthogonal to $\bone_n$. So,  Proposition \ref{prop:singularvalues} implies 
   $$\|\jac(a\score(\alpha), b \score(\beta))\|^2 \leq a^2\|\score(\alpha)\|^2+b^2 \|\score(\beta)\|^2.$$
   Combining this with Observation \ref{obs:scorevecs} yields 
   $$(a+b)^2 \|\score(\gamma)\|^2 \leq a^2\|\score(\alpha)\|^2+ b^2\|\score(\beta)\|^2. $$
   Now, by choosing $a= \frac{1}{\|\score(\alpha)\|^2}$ and $b = \frac{1}{\|\score(\beta)\|^2}$, the above inequality turns into 
   $$\left(\frac{1}{\|\score(\alpha)\|^2} + \frac{1}{\|\score(\beta)\|^2} \right)^2 \|\score(\gamma)\|^2 \leq \frac{1}{\|\score(\alpha)\|^2} + \frac{1}{\|\score(\beta)\|^2}, $$
   which after simple algebraic manipulations can be turned into the inequality claimed in Theorem \ref{thm:stamineq}. 
\end{proof}


\subsection{Understanding $\jac$}
\label{sec:understandingjacob}
Let $(\Omega_{\boxplus_n, 1}, \dots, \Omega_{\boxplus_n, n})$ be the coordinate functions of $\roots$, that is $\gamma_i = \rootsi(\alpha, \beta)$. The starting point of our approach to proving Proposition \ref{prop:singularvalues} is the  observation that the matrix $\jac\jac^*$ is related to the Hessians of the functions $\rootsi$. It will be helpful to introduce the notation
$$\hessconv:= \mathrm{Hess}_{\rootsi}.$$  For this discussion it will prove useful to define the $(2n-2)$-dimensional subspace 
$$\calV = \left\{ (u, v) \in \bR^{n}\times \bR^n : u^* \bone_n = v^*\bone_n=0 \right\}. $$
And, given $w\in \bR^{n}\times \bR^n$ and $f: \bR^{n}\times \bR^n\to \bR^n$ we will use $D_wf$ to denote the directional derivative of $f$ in the direction of $w$, that is $D_w = \sum_i w_i\partial_i$.  
\begin{lemma}[The Hessian of $\roots$]
\label{lem:Hessian}
Using the above notation
    \begin{equation}
    \label{eq:Hessian}
    w^* \jac \jac^* w = w^* \Big(I_{n}\oplus I_n- \sum_{i=1}^n \gamma_i \hessconv \Big)w,\quad \forall w\in \calV. 
      \end{equation}
\end{lemma}

\begin{proof}
Fix $w=(u, v)\in \calV$ and define 
$$\alpha(t):= \alpha+ tu, \quad \beta(t):=\beta+tv,  \quad \text{and}\quad  \gamma(t):=\roots(\alpha(t), \beta(t)),$$ and note that the variance additivity  from Proposition \ref{thm:ffconvolution} \ref{item:additivity}) implies that
$$m_2(\gamma(t))-m_1(\gamma(t))^2 = m_2(\alpha(t))+ m_2(\beta(t))-(m_1(\alpha(t))^2+m_1(\beta(t))^2). $$
Now, the fact that $(u,v)\in \calV$  implies that the means $m_1(\alpha(t))$ and $m_1(\beta(t))$ are a constant function of $t$ and therefore, again by Proposition \ref{thm:ffconvolution} \ref{item:additivity}), the mean $m_1(\gamma(t))$ is also a constant function of $t$. So, differentiating the above equation twice with respect to $t$ we get
\begin{equation}
\label{eq:secderivatives}
 \left.\partial_t^2 m_2(\gamma(t))\right|_{t=0} = \left.\partial_t^2 \big(m_2(\alpha(t))+m_2(\beta(t))\big)\right|_{t=0} .
 \end{equation}
Now we inspect both sides of the above equation. First
\begin{align}
\nonumber \left. n\, \partial_t^2m_2(\gamma(t))\right|_{t=0} & = \sum_{i=1}^n D_w^2(\gamma_i^2) 
\\ \nonumber & = 2\sum_{i=1}^n \left( (D_w\gamma_i)^2+ \gamma_iD_w^2\gamma_i\right)
\\ \label{eq:secondderofgamma} & = 2\left(w^* \jac \jac^* w+ \sum_{i=1}^n\gamma_i w^* \hessconv w \right). 
\end{align}
Second
\begin{align}
\nonumber n\,\partial_t^2(m_2(\alpha(t))+m_2(\beta(t))) & = \partial_t^2((\alpha+tu)^*(\alpha+tu)+(\beta+tv)^*(\beta+tv)) 
 \\ \nonumber & = 2(u^*u +v^*v) 
 \\ \label{eq:secondderofalphanadbeta} & = 2w^* w.
 \end{align}
Finally, plugging (\ref{eq:secondderofgamma}) and (\ref{eq:secondderofalphanadbeta}) back into (\ref{eq:secderivatives}) yields 
$$w^* \jac \jac^* w+ \sum_{i=1}^n\gamma_i w^* \hessconv w = w^* w,$$
which is equivalent to the advertised result. 
\end{proof}


We now apply a result of Bauschke et al. \cite[Corollary 3.3]{bauschke2001hyperbolic}.

\begin{theorem}[Bauschke et al.]
\label{thm:convexity}
    Let $f\in \bR[x_1, \dots, x_m]$ be a hyperbolic polynomial in the direction $w\in \bR^m$ and for every $a\in \bR^m$ let $\lambda_1(a)\geq  \dots\geq  \lambda_m(a)$ be the roots of $g_a(t):=f(a+tw)$.  Then, for every $k=1, \dots, m$, the function $\sigma_k(a) := \sum_{i=1}^k \lambda_i(a)$ is convex in $a$. 
\end{theorem}

In our context this implies the following. 

\begin{corollary}
\label{cor:psdness}
   For any real numbers $c_1\geq \cdots \geq c_n $, the matrix 
   $\sum_{i=1}^n c_i \hessconv$ is PSD. 
\end{corollary}

\begin{proof}
    Define the multivariate polynomial 
    $$f(x, a_1,\dots, a_n, b_1, \dots,  b_n) := \sum_{\pi\in S_n} \prod_{i=1}^n (x-a_i-b_{\pi(i)}).$$
    Since the above polynomial is homogeneous and the finite free convolution preserves real rootedness, $f$ is hyperbolic in the direction $e_1=(1, 0\cdots, 0)$. Now, by Theorem \ref{thm:convexity} the functions
    $$\sigma_k(x, a, b) = \sum_{i=1}^k \lambda_i(x, a, b)$$
    are convex, where $\lambda_1(x, a, b)\geq \cdots \geq \lambda_n(x, a, b)$ denote the roots of $f((x, a, b)+t e_1)$. And, because the $c_i$ are ordered we moreover have that the function 
    $$L(x, a, b):=\sum_{i=1}^n c_i \lambda_i(x,a, b)$$
    is convex, as it can be written as a positive linear combination of the $\sigma_k$. It follows that $\mathrm{Hess}_{L} = \sum_{i=1}^n c_i \mathrm{Hess}_{\lambda_i}$ at any $(x, a, b)$ is PSD. But, on the other hand, when $x=0$, $a=\alpha$ and $b = \beta$, we have that $\mathrm{Hess}_{\lambda_i} = \hessconv$, which in turn gives that $\sum_{i=1}^n c_i \hessconv$ is PSD. 
\end{proof}


We can now complete the proof of Proposition \ref{prop:singularvalues}. 

\begin{proof}[Proof of Proposition \ref{prop:singularvalues}]
  Let $(u, v)\in \calV$. Then 
  $$\|\jac(u, v)\|^2 = (u, v)^* \jac \jac^* (u, v) = \|u\|^2+\|v\|^2-\sum_{i=1}^n \gamma_i (u, v)^* \hessconv (u, v),$$
where the last equality follows from Lemma \ref{lem:Hessian}. Now, applying Corollary \ref{cor:psdness} with $c_i=\gamma_i$ gives that $\sum_{i=1}^n \gamma_i \hessconv$ is PSD, and hence 
$$\sum_{i=1}^n \gamma_i (u, v)^* \hessconv (u, v)\geq 0.$$
The proof follows from putting the two expressions together. 
\end{proof}

\bibliographystyle{alpha}
\bibliography{BibStam}

\end{document}
